\subsection{自由交换群与自由交换幺半群}

集合 \(\mathbf{Z}^{\mathbf{x}}\) （由所有 \(\mathbf{X}\) 到 \(\mathbf{Z}\) 中的映射组成）在由 \((\alpha + \beta)(x) = \alpha(x) + \beta(x)(\alpha, \beta \in \mathbf{Z}^{\mathbf{x}}, x \in \mathbf{X})\) 定义的合成律下构成一个交换群； \(\mathbf{Z}^{\mathbf{x}}\) 的元素有时称为重指标。单位元记作 0，它是恒取值 0 的常值映射。对 \(\alpha \in \mathbf{Z}^{\mathbf{x}}\) ，集合 \(S_{\alpha}\) ，由 \(x \in \mathbf{X}\) （使得 \(\alpha(x) \neq 0\) ）组成，称为 \(\alpha\) 的支集；于是 \(S_0 = \varnothing\) 且 \(S_{\alpha + \beta} \subset S_{\alpha} \cup S_{\beta}\) 对 \(\alpha, \beta\) （在 \(\mathbf{Z}^{\mathbf{x}}\) 中）成立。因此，集合 \(\mathbf{Z}^{(\mathbf{x})}\) （由具有有限支集的映射 \(\alpha: \mathbf{X} \to \mathbf{Z}\) 组成）是 \(\mathbf{Z}^{\mathbf{x}}\) 的一个子群，称为在 \(\mathbf{X}\) 上构造的自由交换群。

对所有 \(x \in \mathbf{X}\) ，设 \(\delta_x\) 表示 \(\mathbf{Z}^{(\mathbf{x})}\) 中由下式定义的元素：

\[
\delta_ {x} (y) = \left\{ \begin{array}{l l} 1 & \text { if } y = x \\ 0 & \text { if } y \neq x. \end{array} \right.\tag{10}
\]

此外，对 \(\alpha \in \mathbf{Z}^{(\mathbf{x})}\) ，整数 \(|\alpha|\) ，即 \(\alpha\) 的长度，由公式定义

\[
| \alpha | = \sum_ {x \in X} \alpha (x).\tag{11}
\]

以下关系立即成立：

\[
\alpha = \sum_ {x \in \mathbf {X}} \alpha (x) \cdot \delta_ {x}\tag{12}
\]

\[
\left| \delta_ {x} \right| = 1, \quad \left| 0 \right| = 0\tag{13}
\]

\[
| \alpha + \beta | = | \alpha | + | \beta |\tag{14}
\]

对 \(\alpha, \beta\) （在 \(\mathbf{Z}^{(\mathbf{X})}\) 中）和 \(x\) （在 \(\mathbf{X}\) 中）成立。

序关系 \(\alpha \leqslant \beta\) 在 \(\mathbf{Z}^{(\mathbf{x})}\) 中由条件 \(\alpha(x) \leqslant \beta(x)\) （对所有 \(x \in \mathbf{X}\) ）定义。关系 \(\alpha \leqslant \beta\) 和 \(\alpha' \leqslant \beta'\) 蕴含 \(\alpha + \alpha' \leqslant \beta + \beta'\) , \(|\alpha| \leqslant |\beta|\) 和 \(-\alpha \geqslant -\beta\) ；进一步，关系 \(\alpha \leqslant \beta\) 等价于 \(\beta - \alpha \geqslant 0\) 。由满足 \(\alpha \geqslant 0\) 的元素（在 \(\mathbf{Z}^{(\mathbf{x})}\) 中）组成的集合记作 \(\mathbf{N}^{(\mathbf{x})}\) ；它是 X 到 N 的具有有限支集的映射组成的集合，并且是 \(\mathbf{Z}^{(\mathbf{x})}\) 的一个子幺半群，称为在 X 上构造的自由交换幺半群。长度为 1 的元素是 \(\mathbf{N}^{(\mathbf{x})} = \{0\}\) 中的极小元素，并且构成 \(\delta_x\) ( \(x \in \mathbf{X}\) ) 的集合。

幺半群 \(\mathbf{N}^{(\mathbf{x})}\) 和群 \(\mathbf{Z}^{(\mathbf{x})}\) 具有以下泛性质。

命题10. 设 M 为交换幺半群（相应地，群），f 为 X 到 M 的映射。存在且仅存在一个从 \(\mathbf{N}^{(\mathbf{x})}\) （相应地 \(\mathbf{Z}^{(\mathbf{x})}\) ）到 M 中的同态，使得 \(g(\delta_x) = f(x)\) 对所有 \(x \in \mathbf{X}\) 成立。若 \(\mathbf{M}\) 以加法记写，则 \(g(\alpha) = \sum_{x \in \mathbf{X}} \alpha(x) \cdot f(x)\) 对所有 \(\alpha\) （在 \(\mathbf{N}^{(\mathbf{x})}\) （相应地 \(\mathbf{Z}^{(\mathbf{x})}\) ）中）成立。

设 \(g\) 是 \(\mathbf{N}^{(\mathbf{X})}\) （相应地 \(\mathbf{Z}^{(\mathbf{X})}\) ）到 \(\mathbf{M}\) 中的同态，使得 \(g(\delta_x) = f(x)\) 对所有 \(x \in \mathbf{X}\) 成立。对所有 \(\alpha\) （在 \(\mathbf{N}^{(\mathbf{X})}\) （相应地 \(\mathbf{Z}^{(\mathbf{X})}\) ）中），由 (12) 可得

\[
g (\alpha) = \sum_ {x \in \mathbf {X}} \alpha (x). g (\delta_ {x}) = \sum_ {x \in \mathbf {X}} \alpha (x). f (x),
\]

由此得出g的唯一性。

对所有 \(\alpha\) （在 \(\mathbf{N}^{(\mathbf{x})}\) （相应地 \(\mathbf{Z}^{(\mathbf{x})}\) ）中），我们写 \(g(\alpha) = \sum_{x\in \mathbf{X}}\alpha (x).f(x)\) 。那么显然 \(g(0) = 0\) ；对 \(\alpha, \beta\) （在 \(\mathbf{N}^{(\mathbf{x})}\) （相应地 \(\mathbf{Z}^{(\mathbf{x})}\) ）中），

\[
\begin{array}{r l} g (\alpha + \beta) & = \sum_ {x \in \mathbf {X}} (\alpha (x) + \beta (x)). f (x) \\ & = \sum_ {x \in \mathbf {X}} [ \alpha (x). f (x) + \beta (x). f (x) ] \\ & = \sum_ {x \in \mathbf {X}} \alpha (x). f (x) + \sum_ {x \in \mathbf {X}} \beta (x). f (x) \\ & = g (\alpha) + g (\beta) \end{array}
\]

，因此 \(g\) 是 \(\mathbf{N}^{(\mathbf{x})}\) （相应地 \(\mathbf{Z}^{(\mathbf{x})}\) ）到 \(\mathbf{M}\) 中的同态。此外，对 \(y\) （属于 \(\mathbf{X}\) ），

\[
g (\delta_ {y}) = \sum_ {x \in \mathbf {X}} \delta_ {y} (x) \cdot f (x);
\]

现在 \(\delta_y(x).f(x) = 0\) 对 \(x\neq y\) 成立，而 \(\delta_y(y).f(y) = f(y)\) ，因此 \(g(\delta_y) = f(y)\) 。

设 \(u: X \to Y\) 是一个映射。根据命题 10，存在且仅存在一个从 \(\mathbf{Z}^{(X)}\) 到 \(\mathbf{Z}^{(Y)}\) 中的同态，它把 \(\delta_x\) 映到 \(\delta_{u(x)}\) （对所有 \(x \in X\) ）。它记作 \(\mathbf{Z}^{(w)}\) ；立即可见它把 \(\alpha \in \mathbf{Z}^{(X)}\) 映到元素 \(\beta \in \mathbf{Z}^{(Y)}\) ，后者由下式定义：

\[
\beta (y) = \sum_ {x \in u ^ {- 1} (y)} \alpha (x).\tag{15}
\]

如同广群的情形（第 1 号），公式 \(\mathbf{Z}^{(v\circ u)} = \mathbf{Z}^{(v)}\circ \mathbf{Z}^{(y)}\) 对每个映射 \(v\colon \mathbf{Y}\to \mathbf{Z}\) 成立；还可证明 \(\mathbf{Z}^{(u)}\) 在 \(u\) 为单射（相应地满射、双射）时是单射（相应地满射、双射）。

设 \(S\) 是 \(X\) 的一个子集；若 \(i\) 是 \(S\) 到 \(X\) 中的单射，则映射 \(f = \mathbf{Z}^{(i)}\) 是 \(\mathbf{Z}^{(S)}\) 到子群 \(H\) （\(\mathbf{Z}^{(X)}\) 中由元素 \(\delta_s\) （对 \(s \in S\) ）生成的）上的同构。由 (15)，

\[
(f (\alpha)) (x) = \left\{ \begin{array}{l l} \alpha (x) & \text { if } x \in S \\ 0 & \text { if } x \in X - S \end{array} \right.
\]

因此 H 是 \(\mathbf{Z}^{(\mathbf{x})}\) 中支集包含于 S 中的元素组成的集合。此后， \(\mathbf{Z}^{(\mathbf{s})}\) 将借助 f 与 H 等同。

公式 (15) 表明， \(\mathbf{Z}^{(u)}\) 在 \(\mathbf{N}^{(\mathbf{X})}\) 上的限制诱导一个同态 \(\mathbf{N}^{(u)}\) （从 \(\mathbf{N}^{(\mathbf{X})}\) 到 \(\mathbf{N}^{(\mathbf{Y})}\) 中）。则 \(\mathbf{N}^{(v\circ u)} = \mathbf{N}^{(v)}\circ \mathbf{N}^{(u)}\) 对每个映射 \(v\colon \mathbf{Y}\to \mathbf{Z}\) 成立；进一步， \(\mathbf{N}^{(u)}\) 在 \(u\) 为单射（相应地满射、双射）时是单射（相应地满射、双射）。如果 \(S\) 是 \(\mathbf{X}\) 的一个子集，

\[
\mathbf {N} ^ {(\mathbf {S})} = \mathbf {Z} ^ {(\mathbf {S})} \cap \mathbf {N} ^ {(\mathbf {X})}.
\]

注：设 M 为严格正整数的乘法幺半群，P 为素数集合（§ 4，第 10 号，定义 15）。由命题 10，存在一个同态 u ，它把 \(\mathbf{N}^{(\mathfrak{P})}\) 映入 M，其特征是 \(u(\delta_{p}) = p\) （对每个素数 p）。则 \(u(\alpha) = \prod_{p \in \mathfrak{P}} p^{\alpha(p)}\) 对 \(\alpha\) （属于 \(\mathbf{N}^{(\mathfrak{P})}\) ）成立，且 § 4 第 10 号的定理 7 表明 u 是 \(\mathbf{N}^{(\mathfrak{P})}\) 到 H 上的一个同构。

